% Chapter 3

\chapter{User Manual}
\label{Chapter3}
\setlength{\parindent}{0pt}

%----------------------------------------------------------------------------------------
\section{Requirements} The project was developed and tested using the following software versions:
\begin{itemize}
	\item \textbf{Geth}: 1.13.15-stable-c5ba367e
	\item \textbf{Node.js}: v22
	\item \textbf{Hardhat}: v3.9.0
	\item \textbf{Solidity}: v0.8.28
	\item \textbf{Apache Maven}: v3.9.9
\end{itemize}

Before executing the project, the Bitcoin blockchain blocks required by the off-chain scanner must be manually copied into the \texttt{bt\_chain\_blocks} directory located in the project root.

%----------------------------------------------------------------------------------------
\section{Project Structure}

The project is organized into independent modules that separate blockchain data, smart contracts, off-chain services, deployment scripts, and utilities.

\begin{itemize}
	\item \textbf{bt\_chain\_blocks/}: local Bitcoin blockchain used by the off-chain scanner.
	\item \textbf{chain/}: local Ethereum configuration (genesis, accounts, and Geth data).
	\item \textbf{hardhat/}: Solidity smart contracts and deployment scripts.
	\item \textbf{offchain-scanner/}: Java application that extracts UTXO data from the Bitcoin blockchain.
	\item \textbf{python/}: Python scripts for deployment, oracle management, automated voting, and demonstration.
	\item \textbf{state/}: stores deployment state shared by the Python components.
	\item \textbf{tool.py}: orchestration script used to manage the entire system.
	\item \textbf{venv/}: Python virtual environment with project dependencies.
\end{itemize}

%----------------------------------------------------------------------------------------
\section{Orchestration Script}

The \texttt{tool.py} script provides an interactive command line interface to automate the setup and execution of the system.
The available operations are:

\begin{enumerate}
	\item Run the Off-chain Scanner.
	\item Create the Python virtual environment and install dependencies.
	\item Initialize the local Ethereum blockchain.
	\item Start the Geth node with mining and exposed HTTP JSON-RPC interface.
	\item Create accounts and deploy contracts.
	\item Start the Oracle Daemon.
	\item Start the Auto Voter.
	\item Run the lending demonstration.
	\item Attach to the running Geth console.
\end{enumerate}

The recommended execution order follows the list above. Long-running services (Geth, Oracle Daemon, and Auto Voter) must remain active while the system is running; therefore, each of them should be started in a separate terminal by executing \texttt{python3 tool.py} and selecting the corresponding option.